\documentclass{article}
\usepackage[utf8]{inputenc}
\usepackage{graphicx}
\usepackage{pgfplots}
\usepackage{amsmath}
\usepackage{amssymb}
\usepackage{amsfonts}
\usepackage{geometry}
\usepackage[thmmarks, amsmath, thref]{ntheorem}
\usepackage{mdframed}
\usepackage{amsthm}
\usepackage{physics}
\usepackage{proof}
\usepackage{derivative}
\usepackage{comment}
\usepackage{hyperref}
\usepackage{wrapfig}
\usepackage{amsthm}
\usepackage{xcolor}
\renewcommand{\theenumi}{\roman{enumi}}

\theoremstyle{nonumberplain}
\theoremheaderfont{\itshape}
\theorembodyfont{\upshape}
\theoremseparator{.}
\theoremsymbol{\ensuremath{\square}}
\theoremsymbol{\ensuremath{\blacksquare}}
\newtheorem{solution}{Solution}
\theoremseparator{. }
\theoremsymbol{\mbox{\texttt{;o)}}}

\pgfplotsset{compat=1.18}
\begin{document}

\textbf{Realiza un trabajo de investigación en el que contestes a todas estas preguntas, hazlo de forma lógica y ordenada, que no sea obvio que estás respondiendo preguntas sino que sea más bien como algo que aprenderías en un libro. El mayor desafío es enlazar todo con lenguaje fluido, no forzado, y poder pasar de una pregunta a otra con transiciones suaves, de modo que no parezcan bloques apilados.}

\begin{enumerate}
    \item[\textbf{I.}] ¿Cómo encontraban los antiguos navegantes el norte usando las estrellas?
    
    \item[\textbf{II.}] ¿Cómo puedo encontrar la Estrella Polar usando ``las punteras'' de la Osa Mayor?\\
    \textit{(Debes de decir también quiénes son las punteras, de preferencia adjuntar una foto editada en el programa de tu preferencia)}
    
    \item[\textbf{III.}] ¿Por qué en algunos lugares del mundo no se ve la Estrella Polar?
    
    \item[\textbf{IV.}] Si salgo al atardecer y veo hacia dónde se pone el Sol, ¿siempre se pondrá exactamente en el mismo lugar?
    
    \item[\textbf{V.}] ¿Cuántos días debería tener un año exactamente si consideramos el tiempo real que tarda la Tierra en dar una vuelta al Sol?\\
    \textit{(¿Por qué existe el año bisiesto?)}
    
    \item[\textbf{VI.}] ¿Por qué tenemos estaciones del año y cómo cambian las posiciones del Sol a lo largo del año?
    
    \item[\textbf{VII.}] ¿Qué ocasiona las olas del mar?
    
    \item[\textbf{VIII.}] ¿Por qué a veces vemos la Luna de día?
    
    \item[\textbf{IX.}] ¿Por qué el cielo es azul y por qué cambia de color al atardecer?
    
    \item[\textbf{X.}] Si veo las estrellas durante la noche, ¿se mueven? ¿Se mueven de verdad o es la Tierra la que se mueve?
    
    \item[\textbf{XI.}] ¿Qué vería un niño que vive en el hemisferio sur cuando mira las estrellas?\\
    \textit{(Comparar con lo que vemos en el hemisferio norte)}
\end{enumerate}

\section*{Rúbrica de Evaluación del Trabajo de Investigación}

\begin{tabular}{|p{9cm}|p{3cm}|}
\hline
\textbf{Criterio} & \textbf{Puntaje} \\
\hline
\textbf{Contenido:} Responde correctamente a todas las preguntas, demostrando comprensión y búsqueda de información. & 4 puntos \\
\hline
\textbf{Redacción:} El trabajo está escrito de manera fluida, con transiciones naturales entre temas, sin que parezca una lista de respuestas. & 3 puntos \\
\hline
\textbf{Creatividad y Presentación:} Uso de imágenes, ejemplos, comparaciones o dibujos que hagan el trabajo más interesante y visual. & 2 puntos \\
\hline
\textbf{Ortografía y claridad:} Buena ortografía, gramática y claridad en las ideas. & 1 punto \\
\hline
\end{tabular}

\vspace{0.5cm}
\newpage
\textbf{Penalizaciones:}

\begin{itemize}
    \item \textbf{-1 punto} si se copia literalmente la pregunta en el texto.
    \item \textbf{-1 punto} si las respuestas están claramente separadas por bloques sin conexión, con transiciones forzadas o poco naturales.
    \item \textbf{-1 punto} si no se entrega imagen o dibujo para el apartado de las punteras de la Osa Mayor (pregunta II).
    \item \textbf{-1 punto} si el trabajo es demasiado corto o incompleto, pero \textbf{-2 puntos} si el trabajo es demasiado largo, debes saber identificar la información importante.
    \item \textbf{-4 puntos} si se nota que solo está transcribiendo cambiando una que otra palabra para su trabajo (es muy facil de notar, sobre todo porque el estilo de escritura aunque cambies palabras al ver la oración total cambiará notoriamente de una sentencia a otra).
\end{itemize}

\vspace{0.5cm}

\textbf{Puntaje máximo: 10 puntos}


\end{document}

